\section{Introduction}
\label{sec:current-statement}

For fixed $1\le a\le2$, we say that \emph{Proposition $(1+a)$} holds
if every sufficiently large even integer $N$ has a representation
\[
 N=p+rq,\qquad r\le q^{a-1},\qquad
 p,q\text{ prime},\qquad r=1\text{ or prime}.
\]
When $a=1$, the condition forces $r=1$, giving the binary Goldbach
conjecture for sufficiently large even integers. When $a=2$, ordering
the two factors gives Chen's theorem. The intermediate parameters ask
how much smaller one prime factor of a composite remainder can be
required to be.

We prove the following theorem.

\begin{theorem}
\label{cor:headline}
Proposition $(1+1.75)$ holds unconditionally. Thus every sufficiently
large even integer can be written as $N=p+rq$, where $p,q$ are primes,
$r$ is either $1$ or a prime, and $r\le q^{3/4}$.
\end{theorem}

The thresholds for $N$ in Theorems~\ref{cor:headline}--\ref{thm:current-main}
and Corollaries~\ref{cor:three-quarters}--\ref{cor:current-explicit} are ineffective.

The parameter $1.75$ can be reduced to $1.74816964$ and further to $1.74816865$,
with different quantitative lower bounds for the number of
representations; see Corollary~\ref{cor:current-explicit}.

Li and Liu~\cite[Theorem~1.1]{LiLiu} prove Proposition $(1+1.9)$
and remark that their parameter can be reduced to $1.894$.
Theorem~\ref{cor:headline} strengthens this factor restriction.
Both statements concern every sufficiently large even integer.

For the quantitative argument put $b=a-1$ and count prime summands
with $r\le q^b$. Prime remainders $r=1$ remain admissible, since
$1\le q^b$. The next theorem gives a lower bound for every fixed
$1/2<b<1$.

Throughout the paper, $N>2$ tends to infinity through the even integers,
unsubscripted logarithms are natural, $\log_N x=\log x/\log N$,
and $\Omega(n)$ counts prime factors with multiplicity. Put
\begin{equation}
 C_2=\prod_{\substack{\ell>2\\\ell\text{ prime}}}
 \left(1-\frac1{(\ell-1)^2}\right),\qquad
 C_N=C_2\prod_{\substack{\ell\mid N,\ \ell>2\\\ell\text{ prime}}}
 \frac{\ell-1}{\ell-2},\qquad
 H_N=C_N\frac{N}{\log^2 N}.
 \label{eq:current-normalization}
\end{equation}
For any positive even integer $M$, $C_M$ denotes the same product with $M$ in place of $N$.
We count distinct prime summands by
\begin{align}
 D_{1,2}(N)&=\#\{p<N:p\text{ prime},\ N-p>1,
                                      \ \Omega(N-p)\le2\},
                                      \label{eq:current-d12}\\
 D_b(N)&=\#\{p<N:p\text{ prime},\ N-p=rq,
       \ q\text{ prime},\ r=1\text{ or prime},\ r\le q^b\},
                                      \label{eq:current-db}\\
 B_b(N)&=\#\{p<N:p\text{ prime},\ N-p=rq,
                \ r,q\text{ prime},\ q^b<r\le q\}.
                                      \label{eq:current-bb}
\end{align}
The ordered factorization $r\le q$ is unique. In particular, a composite
remainder contributes once; a square contributes to $B_b$, whereas a prime
remainder contributes to $D_b$.

\begin{theorem}
\label{thm:current-main}
Let
\begin{equation}
 K=\frac{8452171}{1250000}=6.7617368,
 \qquad c=\frac{49044063}{25000000}=1.96176252.
 \label{eq:current-safe-constants}
\end{equation}
For each fixed $1/2<b<1$,
\begin{equation}
 \liminf_{\substack{N\to\infty\\2\mid N}}
      \frac{D_b(N)}{H_N}\ge c-K\log(1/b).
 \label{eq:current-general-theorem}
\end{equation}
\end{theorem}

Within $1/2<b<1$, this lower bound is positive precisely when
$b>\exp(-c/K)$; see Section~\ref{sec:current-transfer}.

\begin{corollary}
\label{cor:three-quarters}
For every sufficiently large even integer $N$,
\[
                  D_{3/4}(N)\ge0.0165H_N.
\]
\end{corollary}

\begin{corollary}
\label{cor:current-explicit}
Every sufficiently large even integer $N$ satisfies
\begin{align}
 D_{0.74816964}(N)&\ge10^{-5}H_N,
                                  \label{eq:current-main-count}\\
 D_{0.74816865}(N)&\ge10^{-6}H_N,
                                  \label{eq:current-secondary-count}\\
 D_{0.9}(N)&\ge1.24889H_N.
                                  \label{eq:current-common-exponent}
\end{align}
\end{corollary}

The first two inequalities establish Propositions $(1+1.74816964)$
and $(1+1.74816865)$ with respective count coefficients $10^{-5}$ and
$10^{-6}$ in the normalization~\eqref{eq:current-normalization}.
The estimate at $b=0.9$ permits a comparison under the same factor
restriction.

\begin{proof}[Proof of Theorem~\ref{cor:headline}]
Corollary~\ref{cor:three-quarters} gives $D_{3/4}(N)\ge0.0165H_N>0$
for every sufficiently large even $N$. Thus Proposition $(1+1.75)$ holds.
\end{proof}

\subsection{Related work and the method}
In the normalization~\eqref{eq:current-normalization}, the result of
Li and Liu~\cite[Theorem~1.1]{LiLiu} has coefficient $0.0004$ at
$b=0.9$; their stated refinement corresponds to $b=0.894$.
At $b=0.9$, our $D_b(N)$ coincides with Li--Liu's $D_{1,1.9}(N)$:
both include the prime-remainder case $r=1$, count each admissible
prime summand once, and use the same singular-series normalization.
Thus the coefficients refer to the same counting function.
Our coefficient follows from the combination $c-K\log(1/b)$: $c$ is
supplied by the lower Chen-count construction, and $K$ by the positive
upper majorant.
The smaller exponents in Corollary~\ref{cor:current-explicit} retain
an explicit positive count. Their discussion of lower and marked sieve
terms~\cite[Sections~3 and 8.2]{LiLiu} also identifies a constraint
that matters here: an improved upper-sieve distribution level must be
matched to the coefficient class of each term where it is used.

The double sieve transfers information between upper and lower sieve
remainders. We use Wu's formulation~\cite{Wu}, with finite positive
transitions for the two remainder functions.
Li~\cite[Theorem~1.1]{RunboLi} gives coefficient $1.9728$ for the total
Chen count in the same normalization. The coefficient $c=1.96176252$
used here is proved within the moving-residue coefficient class and
positive marked families of the factor-restricted argument. Each lower
transition retains its actual residual level and all prefix conditions.

The distribution estimate must be uniform as the residue varies with
$N$. Maynard's triply well-factorable estimate~\cite[Definition~2 and
Theorem~1.1]{MaynardII} has a constant depending on a fixed residue.
Lichtman's large-residue statement~\cite[Theorem~1.7]{Lichtman} has
level $(5-\vartheta)/8-\epsilon=153/256-\epsilon$ at
$\vartheta=7/32$. Here the residue is comparable with the complete
local product. The reciprocal phase
in~\eqref{eq:current-retained-phase} is retained, since replacing it
by $1$ in this range need not give a vanishing error.
Appendices~\ref{app:medium} and~\ref{app:classical}, together with the
factorization in Appendix~\ref{app:coefficients}, establish the
restricted distribution input at support exponent $0.599$ for the
finite upper and lower coefficient families used in this paper.
At each application we retain the selected factorization, full and
complementary square conditions, and positivity of the marked family.

The proof subtracts the excluded semiprimes from the full Chen count.
For the upper estimate we construct a nonnegative divisor majorant,
compute its density, and prove its distribution estimate for the same
coefficient. For the lower estimate we restrict the two-prime and
four-prime terms of an eleven-term Chen weight in pairs. An injection
preserves the pointwise weight inequality, while a lower bound for
Buchstab's function on the restricted fibres improves the coefficient.

The nonnegative upper majorant and its coefficient realization are
constructed here. For the lower bound, we use Wu's double-sieve framework
with the paired deletion and coefficient realizations established here.
Together with the finite transitions and directed enclosures, these
constructions determine the constants $K,c$. Their mathematical
definitions and certificate correspondence are given in
Appendix~\ref{app:numeric}.

\section{Reduction and outline of the proof}
\label{sec:outline}

\paragraph{Proof outline and verification.}
On a first reading, Lemma~\ref{lem:numerical-endpoints} may be taken as
the finite numerical input: follow the upper construction, the eleven-term
lower construction and paired deletion, and then the final subtraction in
Section~\ref{sec:current-transfer}. For a complete verification, return to
the appendices for the distribution estimates, coefficient realizations,
finite transitions and directed enclosures, and check their correspondence
with the Numerical Supplement described in Appendix~\ref{app:numeric}.

Unique factorization gives the exact identity
\begin{equation}
                 D_b(N)=D_{1,2}(N)-B_b(N).
                 \label{eq:reduction}
\end{equation}
Thus it suffices to prove
\begin{equation}
 \liminf_{\substack{N\to\infty\\2\mid N}}
       \frac{D_{1,2}(N)}{H_N}\ge c,
 \qquad
 \limsup_{\substack{N\to\infty\\2\mid N}}
       \frac{B_b(N)}{H_N}\le K\log(1/b).
 \label{eq:two-estimates}
\end{equation}
The factor $\log(1/b)$ comes from the excluded factor region. For
$\beta=b/(1+b)$, the limiting logarithmic region of the ordered bad
semiprimes is $\beta<t\le1/2$, and
\[
 \int_\beta^{1/2}\frac{dt}{t(1-t)}=\log(1/b).
\]
To turn this mass into a count with a prime complement, we apply a
positive upper sieve whose distribution and density have both been
established for its actual coefficients.

Section~\ref{sec:current-distribution}
establishes the distribution estimate for the actual upper and lower
coefficients. Section~\ref{sec:current-upper} constructs the positive
upper sieve and obtains $K$. Sections~\ref{sec:current-chen-weight}
and~\ref{sec:current-lower} give the lower construction; its switched
costs use the upper bound already proved. Section~\ref{sec:current-numeric}
states the finite numerical inequalities and their error allowances.
Finally Section~\ref{sec:current-transfer} proves the second estimate
in~\eqref{eq:two-estimates} and completes the theorem. In particular,
the construction of $K$ uses neither the lower double-sieve tables nor
the Chen count obtained from them.

\paragraph{Choice of parameters.}
At each fixed accuracy, the finite collection of policies and marked queries,
the number of generations, cutoffs defining fixed compact parameter domains,
approximation tolerances and
support reserves are chosen before $N$ tends to infinity. Prime intervals,
weights and arithmetic coefficients may still depend on $N$ through these
fixed rules. The common choices and error propagation are specified in
Section~\ref{sec:lower-limits} and \eqref{coef:parameter-order}; the strict
scalar surpluses in Section~\ref{sec:current-numeric} allow one fixed accuracy
to be chosen for the final count.

\paragraph{Notation index.}
The following symbols keep the same meaning in the main text and appendices.
\begin{center}
\small
\begin{tabular}{@{}>{\raggedright\arraybackslash}p{.23\textwidth}>{\raggedright\arraybackslash}p{.49\textwidth}>{\raggedright\arraybackslash}p{.22\textwidth}@{}}
\toprule Symbol & Meaning & Location\\
\midrule
$E_{\rm up}$; $R_{\rm low}$ & Upper and lower scalar endpoints & \eqref{eq:current-E}; \eqref{eq:current-R}\\[2pt]
$\mathscr E_0,\mathscr E_1,\mathscr E$ & Paired rectangle families and their union & Section~\ref{sec:current-chen-weight}\\[2pt]
$a_{\rm s},h_{\rm s},d_{\rm s}$ & Fixed sieve cutoffs & Section~\ref{sec:current-chen-weight}\\[2pt]
$H(s),h(s)$; $H_m,h_m$ & Upper and lower defects; generation-$m$ envelopes & Section~\ref{sec:current-lower}\\[2pt]
$B$; $\mathscr B$ & Arithmetic prime prefix; geometric rectangle or box & Sections~\ref{sec:current-distribution}, \ref{sec:current-lower}\\[2pt]
$H_N$; $B_{\rm low}$ & Count normalization; each lower error allowance & \eqref{eq:current-normalization}; \eqref{eq:current-B}\\[2pt]
$A_+,A_-$ & Normalized linear-sieve numerators $sF/(2e^\gamma),sf/(2e^\gamma)$ & Section~\ref{sec:current-lower}\\[2pt]
$\Omega$; $\omega_{\rm env}$ & Prime-factor count with multiplicity; Buchstab upper envelope & Section~\ref{sec:current-statement}; \eqref{sv:omega-envelope}\\[2pt]
$\rho_N^{\rm up},\rho_N^{\rm low}$ & Actual upper and lower divisor coefficients; $\rho_N$ denotes a generic admitted coefficient & Sections~\ref{sec:current-upper}, \ref{sec:current-lower}\\[2pt]
\bottomrule
\end{tabular}
\end{center}

We use the standard linear sieve functions $F,f$, normalized by
\[
 sF(s)=2e^\gamma\quad(1\le s\le3),\qquad
 sf(s)=0\quad(1\le s\le2),
\]
and subsequently by
$(sF(s))'=f(s-1)$ and $(sf(s))'=F(s-1)$ on their respective
ranges. Buchstab's function satisfies
\[
 \omega(s)=1/s\quad(1\le s\le2),\qquad
 (s\omega(s))'=\omega(s-1)\quad(s>2).
\]
Here $P^-(n)$ is the least prime factor for $n\ge2$, and $P^-(1)=+\infty$. The finite interval data used for these
functions are discussed in Appendix~\ref{app:numeric}.

\subsection{A positive sieve coefficient}
The upper estimate uses the following single coefficient theorem.
Its three conclusions concern the same divisor coefficient.
Here ``positive'' means that the resulting divisor sum $R_N(n)$ is
a nonnegative majorant. The individual coefficients $\rho_N^{\rm up}(q)$
may have either sign.

\begin{theorem}
\label{thm:positive-coefficient}
Fix $a_*>0$, a bound on the number of labelled prime factors, and a
finite collection of smooth weights $W_N(p_1,\ldots,p_k)$ supported on
\[
 p_i\ge c_*N^{a_*},\qquad p_1\cdots p_k\le C_*N,
\]
with fixed constants $c_*,C_*>0$, uniformly bounded size, and fixed
polylogarithmic bounds for their mixed Euler derivatives. Apart from
these smooth geometric conditions, impose no additional arithmetic
filter. Put $\rho_0=19/96$, $\theta_*=599/1000$, and
\[
 P_N(y)=\prod_{\substack{p<y\\p\nmid N}}p,\qquad
 X_W^{(q)}=\sum_{(p_1\cdots p_k,q)=1}W_N(p_1,\ldots,p_k).
\]
There are real coefficients $\rho_N^{\rm up}(q)$ and a constant $M$, independent
of $N$, with squarefree support $q\le N^{\theta_*}$, $(q,N)=1$,
and $|\rho_N^{\rm up}(q)|\le M$, such that
\begin{equation}
 R_N(n):=\sum_{q\mid n}\rho_N^{\rm up}(q)
   \ge\one_{(n,P_N(N^{\rho_0}))=1}\ge0
                           \qquad(n\in\mathbb Z),\label{eq:core-positivity}
\end{equation}
\begin{equation}
 \sum_q\frac{\rho_N^{\rm up}(q)}{\varphi(q)}
   \le(K+o(1))\frac{C_N}{\log N},\label{eq:core-density}
\end{equation}
\begin{multline}
 \sum_q\rho_N^{\rm up}(q)\left\{
 \sum_{p_1\cdots p_k\equiv N\ (q)}W_N(p_1,\ldots,p_k)
       -\frac{X_W^{(q)}}{\varphi(q)}\right\}\\
   =O_{A,W}(N(\log N)^{-A})\qquad(A>0).
                           \label{eq:core-distribution}
\end{multline}
The finite sum defines $R_N(0)$. The finite policy, approximation
parameters and support margins are fixed before $N$ tends to infinity
through all even integers. The constants may depend on this fixed
collection and its derivative bounds, but are uniform in $N$.
\end{theorem}

The coefficient is constructed in Section~\ref{sec:current-upper}.
The distribution statement is proved first in
Section~\ref{sec:current-distribution} for the precise coefficient
representations needed by that construction.
